\documentclass[12 pt]{article} 
\usepackage{float, amsfonts, amssymb, amsmath, blindtext, multicol, graphicx,
enumitem, listings, hyperref}
\usepackage[font=small,labelfont=bf]{caption}
\graphicspath{ {./images/} }

\oddsidemargin=-0.5cm
\setlength{\textwidth}{6.5in}
\addtolength{\voffset}{-20pt}
\addtolength{\headsep}{25pt} 



\pagestyle{myheadings}                           	
\markright{Anthony Nguyen, Rowan Lochrin\hfill \today \hfill} 	    


\newcommand{\eqn}[0]{\begin{array}{rcl}}
\newcommand{\eqnend}[0]{\end{array} } 

\newcommand{\qed}[0]{$\square$}

\begin{document}
\noindent \textbf{CSC-696D, Takanori Fujiwara} \par
\noindent \textbf{Final project proposal}. \\
\noindent \textbf{Novel visualization techniques for git data}.
\bigskip
\subsection{Problem description}
In this project we will create several novel visual analytic systems for git
data. The tools we develop will help developers quickly gain an understanding of
git repositories. In particular, we seek to give users an understanding of how a
repo changes over time. We will focus on a few key questions.
\begin{enumerate}
    \item Who are the main contributors to a repository?
    \item What were the development priorities for a repository at any given
        time?
    \item Can we identify a time when AI-assisted development became popular on
        a given repository?
    \item What is the best visualization for giving developers an understanding of the
        evolution of project structure over time?
\end{enumerate}
\subsection{Motivation}
Developers often work with large amounts of code and are forced to quickly learn
the customs and conventions of many repositories, this process can be very
time-consuming. By giving developers tools for quickly visualizing the history
of a repository, we can immediately give them an understanding of conventions, key
contributors, and etiquette. \par
Many visualizations have been developed over the years to work with git data,
in this project we want to develop a handful of new visualizations by first
studying visualizations that others have developed and refining them to improve
their visual clarity, augmenting them with additional data or even combining two
visualization techniques into one. \par
Currently we have planned 3 visualizations but may add more or modify some of
them over the course of the project.
\begin{enumerate}
    \item A `stabilized' Sediment Graph based on the work of
        \href{https://github.com/erikbern/git-of-theseus}{Erik Bernhardsson}.
        Bernhardsson proposes a unique way to understand the age of the contents of a
        repo, we seek to improve his designs by evenly spacing out various
        layers of the graph (instead of stacking them), to reduce visual
        distortions from oscillating streams of data.
    \item A commit graph and word cloud style view that allows users to select
        many commits via a time window selection and see a word cloud of
        the commit messages. This will allow users to quickly see the
        development focuses of a repo at a particular time window. This is
        inspired by a view from the paper \textit{Visual Analytics for Understanding Software Development History Through
        Git Metadata Analysis}. We seek to enhance their view by adding a time
        window selection allowing these clouds to be generated for any section
        of the commit history. Another enhancement we could make is creating an
        animation going through the commit history with a sliding window to see how the word
        cloud changes over time.
    \item A grouping-based visualization that clusters files or directories that
        are often modified together, this will help users understand which parts
        of the codebase are most tightly coupled. We seek to animate this view
        to give users an understanding of how that changes over time.
\end{enumerate}
\subsection{Related research papers}
\begin{itemize}
\item
    \href{https://ieeexplore.ieee.org/stamp/stamp.jsp?tp=&arnumber=9222261}{Githru\:
        Visual Analytics for Understanding Software Development History Through
        Git Metadata Analysis} ---
    This work shows an all-in-one visual analytic tool for examining git history
        and provides a lot of great background on git. Many techniques here may
        be useful; in particular, the paper provides a lot of detail on the
        technical aspects of creating visualizations with git metadata. We
        also took inspiration from this paper for our second proposed
        visualization as discussed above.
\item \href{https://mircealungu.com/docs/assets/papers/22-Git-Truck.pdf}{Git-Truck: Hierarchy-Oriented Visualization of
Git Repository Evolution} --- Git-Truck is a hierarchical file-organization-oriented
    visualization tool that represents git repositories as different treemaps.
    Each treemap has its nodes sized by the number of lines of code in the repository, 
    and the color of its nodes can represent a variety of interesting factors
    such as number of commits, top contributors, or even the last change day.
    This is showing the current structure of a repository, rather than the temporal 
    focus we have in this project. Mainly, we emphasize the importance of seeing
    the historical development of features in a repository.
\item \href{https://arxiv.org/pdf/2508.12649}{ChangePrism: Visualizing the
    Essence of Code Changes} --- This paper gives another approach to visualizing
        code changes that relies on semantic analysis. Though this paper's focus
        is more on visualizing individual commits and discrete changes rather
        than the type of broad analysis of larger repositories we hope to do,
        it gives some good background on existing visual analytic techniques.
\end{itemize}
\subsection{Planned technical approach}
Our visualizations will use the \texttt{.git} files for the repos we are studying as a
standardized source of data. We plan to implement our visualizations in
JavaScript and may pull in dependencies for creating charts as needed.
Currently, we do not have our exact technical stack, but our project will run on
Node and use NPM for dependency management.
React is an option, but we're also investigating if a simpler UI framework would
meet our needs. AI use is planned to speed up implementation, models used will
probably be Claude Haiku/Sonnet/Opus 4.x.

\subsection{Selected dataset(s)}
We will focus on three categories of repos in this project.
\begin{itemize}
    \item Large repositories with many contributors like the \href{https://github.com/torvalds/linux}{Linux
        kernel}
    \item Large repositories with few contributors like the \href{https://github.com/emacs-mirror/emacs}{emacs}
    \item Small repositories with few contributors like \href{https://gitlab.com/distant-horizons-team/distant-horizons}{mods} and web
        applications created by small teams or individuals. For fun, we will even feed
        the git repository for this project into our system.
\end{itemize}
\subsection{Expected results and evaluation methodology}
In this project, we seek to create a web application to demonstrate and explain
our visual analytic systems. It will feature the visualizations we've created, 
a short text description of the visualizations—explaining works inspiring the
visualization, the types of analytical questions the visualization seeks to
answer and the strengths and weaknesses of our approach. The app will also feature
a selector that will allow you to view any git repository with any of the
aforementioned visualization techniques. We will then gather feedback from
professional software engineers on the usefulness of our visualizations and
applicability of our tools to their current development workflow.

% evaluation examples include: case studies, user studies, quantitative evaluation



\end{document}
